% Plantilla editable inspirada en el cartel original del Seminario Chichimeca
% COMPILAR CON XELATEX (LaTeX Workshop en VS Code)
\documentclass[final]{beamer}
\usepackage[orientation=portrait,size=a2,scale=1.15]{beamerposter}
\usetheme{gemini}
\usecolortheme{nott}
\usepackage{fontspec}
\usepackage{graphicx}
\usepackage{amsmath,amssymb,mathtools}
\usepackage{tikz}
\usepackage{tikz-cd}
\usetikzlibrary{arrows.meta,positioning,decorations.pathmorphing}
\usepackage{ragged2e}

% ==========================================
% Bibliografía
% ==========================================
\usepackage[
    backend=biber,
    style=numeric,
    sorting=none,
    maxbibnames=2,
    doi=false,
    isbn=false,
    url=false
]{biblatex}

\addbibresource{research2.bib}

% ========= DATOS EDITABLES =========
\title{Propiedades de separación inducidas\\por cocientes compactos}
\author{Juan Carlos Monter Cortés\\ Dr. Luis Ángel Záldivar Corichi}
\institute{Universidad de Guadalajara}
\footercontent{Doctorado en Ciencias en Matemáticas \hfill Seminario Chichimeca de Álgebra \hfill juan.monter2902@alumnos.udg.mx}
\logoleft{\hspace{8ex}\includegraphics[height=5.5cm]{logos/ppgca-logo.png}}
\logoright{\hspace{8ex}\includegraphics[height=5.5cm]{logos/utfpr-logo.png}}

% ========= COLUMNAS =========
\newlength{\sepwidth}
\newlength{\colwidth}
\setlength{\sepwidth}{0.025\paperwidth}
\setlength{\colwidth}{0.45\paperwidth}
\newcommand{\separatorcolumn}{\begin{column}{\sepwidth}\end{column}}

% ========= NOTACION =========
\DeclareMathOperator{\Frm}{Frm}
\DeclareMathOperator{\pt}{pt}
\DeclareMathOperator{\Top}{Top}
\newcommand{\separacion}[1]{\textcolor{red!70!black}{\textbf{#1}}}

\begin{document}
\begin{frame}[fragile,t]
\begin{columns}[t]
\separatorcolumn
\begin{column}{\colwidth}

\begin{block}{Introducción y motivación}
El objetivo de este cartel es presentar nuevas propiedades de separación en la categoría de marcos a 
través de sus cocientes compactos.\\

Una vez enunciadas dichas propiedades, se comparan con algunas propiedades de 
separación ya existentes en en la literatura sobre teoría de marcos. 
\end{block}

\begin{block}{Preliminares}
Un \emph{\textbf{marco}} es un retícula completa $(A,\leq, \wedge, \bigvee, 1, 0)$ que satisface
\[
 a\wedge\bigvee S=\bigvee\{a\wedge s\mid s\in S\}.
\]
Junto con los morfismos de marcos, forman la categoría $\Frm$.
Recordemos que existe una adjunción entre la categoría $\Frm$ y la categoría $\Top$ de espacios topológicos
\[
\begin{tikzcd}
\Frm \arrow[r, "\pt", bend left] & \Top \arrow[l, "\mathcal{O}", bend left]
\end{tikzcd}
\]
Los marco $A$ y $\mathcal{O}(\pt A)$ los relacionamos a traves del morfismo de marcos
\[
\mathcal{U}\colon A\to \mathcal{O}(\pt A),\qquad a\mapsto \mathcal{U}(a)=\{p\in \pt A\mid p(a)=1\}.
\]
En $\Frm$, los cocientes pueden trabajarse a traves de los \emph{\textbf{núcleos}}, 
los cuales son funciones $j\colon A\to A$ que satisfacen:
\[
a\leq j(a),\quad j(a\wedge b)=j(a)\wedge j(b),\quad j^2=j, \qquad (a, b\in A).
\]
De manera particular, si $a\in A$, definimos los siguientes núcleos:
\[
u_a(x)=a\vee x,\qquad v_a(x)=(a\succ x)
\]
Si $j$ es un núcleo, podemos definir los siguientes objetos:
\begin{itemize}
    \item El \emph{\textbf{marco cociente:}} $A_j=\{a\in A\mid j(a)=a\}$.
    \item El \emph{\textbf{filtro de admisibilidad:}} $\nabla(j)=\{a\in A\mid j(a)=1\}$.
\end{itemize}
\textit{\textbf{Teorema:} Para $j$ un núcleo en $A$, el cociente $A_j$ es compacto si y sólo si $\nabla(j)$ es un filtro abierto.}

Si $F\subseteq A$ es un filtro abierto, podemos construir su \emph{\textbf{intervalo de admisibilidad}}, esto es el conjunto de todos 
los núcleos $j$ tales que $\nabla(j)=F$. Dicho intervalo lo denotamos por 
\(
[v_F, w_F].
\)

En la categoría $\Frm$, se tiene la siguiente relación entre las distintas propiedades de separación:
\[\begin{tikzcd}
	&& {\mathbf{(reg)}} && \\
	{\mathbf{(fit)}} && {\mathbf{(sH)}} && {\mathbf{(H)}+\mathbf{(sfit)}} \\
	{\mathbf{(sfit)}} && {T_1} && {\mathbf{(H)}}
	\arrow[Rightarrow, from=1-3, to=2-1]
	\arrow[Rightarrow, from=1-3, to=2-3]
	\arrow[Rightarrow, from=1-3, to=2-5]
	\arrow[Rightarrow, from=2-1, to=3-1]
	\arrow[squiggly, no head, from=2-1, to=3-3]
	\arrow[squiggly, no head, from=2-1, to=3-5]
	\arrow[squiggly, no head, from=2-3, to=2-1]
	\arrow[squiggly, no head, from=2-3, to=2-5]
	\arrow[squiggly, no head, from=2-3, to=3-1]
	\arrow[Rightarrow, from=2-3, to=3-3]
	\arrow[Rightarrow, from=2-3, to=3-5]
	\arrow[Rightarrow, from=2-5, to=3-1]
	\arrow[squiggly, no head, from=2-5, to=3-3]
	\arrow[Rightarrow, from=2-5, to=3-5]
	\arrow[squiggly, no head, from=3-1, to=3-3]
	\arrow[Rightarrow, from=3-5, to=3-3]
\end{tikzcd}\]
Información sobre el diagrama y teoría de marcos puede ser consultada en \cite{Johnstone82, PicadoPultr, picado2021separation}.
\end{block}

\begin{alertblock}{Idea principal}
Siguiendo las ideas de Sexton y Simmons presentadas en \cite{sexton2006ordinal,sexton2006point}, definimos 4 clases de marcos relacionados con núcleos 
que producen cocientes compactos. 

Por el Teorema de Hofmann-Mislove, existe una correspondencia biyectiva entre 
\[
F\in A^\wedge\longleftrightarrow Q\in \mathcal{Q}S\qquad (S=\pt A).
\]
Cada par relacionado por dicha correspondencia lo llamamos \emph{\textbf{par HM}}.

Considerando $\mathcal{O}S\in \Frm$ y el conjunto $\mathcal{O}S^\wedge$, tenemos el par HM $(G\in \mathcal{O}S^\wedge, Q\in \mathcal{Q}S)$, donde $G\in \mathcal{O}S^\wedge$. 
Así, tenemos una terna HM $(F, Q, G)$ a la cual le podemos asociar los intervalos de admisibilidad $[v_F, w_F]$ y $[v_G, w_G]$.

Si $Q\in \mathcal{Q}S$, definimos el núcleo $[Q']\colon \mathcal{O}S\to \mathcal{O}S$ como
\[
[Q'](U)=(Q'\cup U)^\circ\quad \text{y}\quad v_G\leq [Q']\leq w_G.
\]
\textit{\textbf{Definición:} Sea $A$ un marco y $(F, Q)$ un par HM y $[v_F, w_F]$ su respectivo intervalo de admisibilidad. Decimos que $A$ es:
\begin{itemize}
\item \separacion{eficiente}, si $v_F=u_d$, donde $d=v_F(0)$, para cada filtro abierto $F$.
\item \separacion{KC}, si para $j\in [v_F, w_F]$, $j=u_a$, para algún $a\in A$ y todo filtro abierto $F$.
\item \separacion{apilado}, si $v_F(0)=(\mathcal{U}_*\circ [Q']\circ \mathcal{U})(0)$, para cada par HM $(F, Q)$.
\item \separacion{fuertemente apilado}, si $v_F=\mathcal{U}_*\circ [Q']\circ \mathcal{U}$, para cada par HM $(F, Q)$.
\item \separacion{débilmente apilado}, si $v_F=\mathcal{U}_*\circ v_G\circ \mathcal{U}$, para cada terna HM $(F, Q, G)$.
\end{itemize}}

Se puede verificar que estos marcos se relacionan de la siguiente manera:
\[\begin{tikzcd}
	{\text{KC}} & {\text{Eft.}} && {\text{Ap.}} & {\text{F.Ap.}} & {\text{D.Ap.}}
	\arrow[Rightarrow, from=1-1, to=1-2]
	\arrow[Rightarrow, from=1-5, to=1-4]
	\arrow[Rightarrow, from=1-5, to=1-6]
\end{tikzcd}\]
\end{alertblock}

%\begin{block}{Ejemplos}
%Aquí puedes incorporar ejemplos, contraejemplos y cálculos de núcleos.
%\[
% u_a(x)=a\vee x.
%\]
%\end{block}
\end{column}
\separatorcolumn
\begin{column}{\colwidth}

\begin{block}{Herencia, preservación y conservatividad}
Las clases de marcos definidos anteriormente satisfacen las siguintes propiedades.

\separacion{Teorema:} Los marcos KC, eficientes y fuertemente apilados son hereditarios.

\separacion{Teorema:} Los marcos eficientes se preservan bajo coproductos arbitrarios. Los marcos KC no se preservan
bajo coproductos.

\separacion{Teorema:} Las nociones de apilamiento son conservativas.

\end{block}

\begin{alertblock}{Relación con axiomas de separación}
\separacion{Proposición:} Las siguientes afirmaciones son ciertas:
\begin{enumerate}
    \item Existen marcos fuertemente apilados que no son $T_1$
    \item Existen marcos $T_1$ y fuertemente apilados que no son espaciales.
    \item Existen marcos $\mathbf{(H)}$ que no son fuertemente apilados.
    \item Existen marcos KC compactos que no son $\mathbf{(reg)}$.
    \item Existen marcos KC que no son $\mathbf{(sfit)}$.
    \item $\mathbf{(sH)}$ implica KC.
\end{enumerate}
\end{alertblock}

\begin{block}{Interacción entre las propiedades inducidas por cocientes compactos}
\separacion{Proposición:} Las siguientes afirmaciones son ciertas:
\begin{enumerate}
    \item Si $A\cong \mathcal{O}S$, entonces KC $\iff$ Eficiente. Lo anterior no es cierto
    para marcos arbitrarios.
    \item Si $A\in \Frm$ es eficiente $\mathbf{(fit)}$, entonces $A$ es KC.
    \item Si $A\in \Frm$ es eficiente fuertemente apilado, entonces $A$ es KC.
    \item Si $A\in \Frm$ es KC, entonces fuertemente apilado $\iff$ apilado.
    \item Si $A\in \Frm$ es eficiente débilmente apilado, entonces $A$ es fuertemente apilado.
    \item Si $A\in \Frm$ es KC $\mathbf{(sfit)}$, entonces $A$ es fuertemente apilado.
    \item Existen marcos KC que no son fuertemente apilados. 
\end{enumerate}
\end{block}

\begin{exampleblock}{Conclusiones}
Las propiedades definidas a través de los cocientes compactos tienen relación con la teoría
de separación en $\Frm$. Lo anterior puede permitir estudiar construcciones como la del marco
de parches a través de dichas deficiones.

Los siguientes diagramas muestran un resumen de la interacción entre algunas propiedades de 
separación y las nuevas propiedades definidas.
\[
\begin{tikzcd}
	&&& {\mathbf{(reg)}} &&& \\
	{\mathbf{(fit)}} &&& {\mathbf{(sH)}} & {\mathbf{(H)}} && {\mathbf{(sH)}+\mathbf{(sfit)}} \\
	&&& {\text{KC}} &&& {\text{KC}+\mathbf{(sfit)}} \\
	{\mathbf{(sfit)}} &&& {\text{Eft}} &&& {\text{F.Ap.}} \\
	&&& {T_1} &&& {\text{Ap.}}
	\arrow[Rightarrow, from=1-4, to=2-1]
	\arrow[Rightarrow, from=1-4, to=2-4]
	\arrow[Rightarrow, from=1-4, to=2-7]
	\arrow[squiggly, no head, from=2-1, to=2-4]
	\arrow[squiggly, no head, from=2-1, to=3-4]
	\arrow[Rightarrow, from=2-1, to=4-1]
	\arrow[Rightarrow, from=2-1, to=5-4]
	\arrow[Rightarrow, from=2-4, to=2-5]
	\arrow[Rightarrow, from=2-4, to=3-4]
	\arrow[squiggly, no head, from=2-5, to=4-7]
	\arrow[Rightarrow, from=2-5, to=5-4]
	\arrow[squiggly, no head, from=2-5, to=5-7]
	\arrow[Rightarrow, from=2-7, to=2-5]
	\arrow[Rightarrow, from=2-7, to=3-4]
	\arrow[Rightarrow, from=2-7, to=3-7]
	\arrow[Rightarrow, from=3-4, to=4-4]
	\arrow[squiggly, no head, from=3-4, to=4-7]
	\arrow[squiggly, no head, from=3-4, to=5-7]
	\arrow[Rightarrow, from=3-7, to=4-4]
	\arrow[Rightarrow, from=3-7, to=4-7]
	\arrow[squiggly, no head, from=4-1, to=2-4]
	\arrow[squiggly, no head, from=4-1, to=3-4]
	\arrow[squiggly, no head, from=4-1, to=4-4]
	\arrow[squiggly, no head, from=4-1, to=5-4]
	\arrow[squiggly, no head, from=4-4, to=4-7]
	\arrow[Rightarrow, from=4-4, to=5-4]
	\arrow[squiggly, no head, from=4-7, to=5-4]
	\arrow[Rightarrow, from=4-7, to=5-7]
	\arrow[squiggly, no head, from=5-7, to=5-4]
\end{tikzcd}
\]
\[\begin{tikzcd}
	{\text{Eft}+\mathbf{(fit)}} & {\text{KC}} & {\text{Eft.}+\text{F.Ap.}} \\
	{\text{KC}+\mathbf{(sfit)}} & {\text{F. Ap.}} & {\text{Eft.}+\text{D.Ap.}}
	\arrow[Rightarrow, from=1-1, to=1-2]
	\arrow[Rightarrow, from=1-3, to=1-2]
	\arrow[Rightarrow, from=2-2, to=2-1]
	\arrow[Rightarrow, from=2-3, to=2-2]
\end{tikzcd}\]
\end{exampleblock}

\begin{block}{Referencias}
    \begingroup
    \footnotesize
    \printbibliography[heading=none]
    \endgroup
\end{block}

\end{column}
\separatorcolumn
\end{columns}
\end{frame}
\end{document}
